\documentclass[11pt]{article}
\usepackage[T1]{fontenc}
\usepackage[utf8]{inputenc}
\usepackage[english]{babel}
\usepackage{lmodern}
\usepackage{amsmath}
\usepackage{amssymb}
\usepackage{graphicx}
\usepackage{microtype}
\setlength{\emergencystretch}{4em}

\begin{document}

\section{Implemented spectral calibration model}

\subsection{Indices, data, and campaign preprocessing}

Let $s\in\{1,\dots,M\}$ index sensors, let $k\in\{1,\dots,K\}$ index aligned calibration experiments, and let $j\in\{1,\dots,J\}$ index the discrete frequency bins of the common calibration band. After preprocessing, the calibration data are represented by an observation tensor
\[
Y_{s,k,j}>0,
\qquad
(s,k,j)\in\{1,\dots,M\}\times\{1,\dots,K\}\times\{1,\dots,J\},
\]
in linear power units.

Before alignment, sensor $s$ provides a timestamped series
\[
\{(\tau_{s,r},x_{s,r,1},\dots,x_{s,r,J})\}_{r=1}^{R_s},
\]
where $\tau_{s,r}$ is the record timestamp and the $x_{s,r,j}$ are PSD values
in dB.

The implemented campaign preprocessing stage does the following.

\begin{enumerate}
    \item It loads one campaign directory from the \texttt{data/campaigns/} root and parses one CSV per sensor.
    \item It sorts each sensor series by timestamp and validates that every retained sensor already uses the same PSD frequency grid.
    \item It chooses the shortest sensor series as the alignment anchor, so every retained realization is observable across the full sensor set.
    \item It aligns records greedily by timestamp: for each anchor timestamp, every other sensor contributes the nearest later record within an alignment tolerance, while preserving time order and dropping unmatched rows.
    \item It converts the aligned PSD tensor from dB to linear power, producing $Y_{s,k,j}$.
    \item It computes ranking diagnostics on the aligned dB tensor and chooses a reliable sensor $s_\star$ as the highest record-wise correlation winner that is \emph{not} flagged as a campaign-wide PSD-distribution outlier; if every retained sensor is an outlier, it falls back to the record-wise winner.
    \item It sets the nominal gain baseline to unity for every sensor and frequency bin.
\end{enumerate}

The alignment tolerance is either supplied explicitly by the caller or inferred
from the median campaign sampling period and then clipped to a conservative
range. The aligned experiment timestamp stored by the adapter is the median
timestamp across sensors for each retained realization.

Denote the resulting nominal gain baseline by $G^{(0)}_{s,j}>0$. In the active
campaign workflow the baseline is
\begin{equation}
G^{(0)}_{s,j}=1,
\qquad
s=1,\dots,M,\ j=1,\dots,J.
\label{eq:nominal-baseline}
\end{equation}

So the learned calibration is a \emph{relative} network calibration around a
unit reference scale. Unlike older versions of the project, the current
repository does not ingest external per-sensor nominal response CSV files.

\subsection{Observation model}

The fitted calibration model is
\begin{equation}
Y_{s,k,j}
=
G_{s,j}\,S_{k,j}
+
N_{s,j}
+
\varepsilon_{s,k,j},
\label{eq:model}
\end{equation}
where:
\begin{itemize}
    \item $S_{k,j}\ge 0$ is the latent common spectrum at experiment $k$ and bin $j$;
    \item $G_{s,j}>0$ is the multiplicative gain of sensor $s$;
    \item $N_{s,j}>0$ is the additive sensor floor;
    \item $\varepsilon_{s,k,j}$ is a residual term.
\end{itemize}

Conditionally on the current latent spectra, the working likelihood is Gaussian:
\begin{equation}
\varepsilon_{s,k,j}\mid S_{k,j}\sim \mathcal N(0,\sigma^2_{s,j}),
\label{eq:gaussian}
\end{equation}
with sensor- and frequency-dependent residual variance $\sigma^2_{s,j}>0$.

\subsection{Reparameterization used by the implementation}

The code does not optimize directly in $(G_{s,j},N_{s,j})$. Instead it uses the parameterization
\begin{equation}
G_{s,j}
=
G^{(0)}_{s,j}\exp(\gamma_{s,j}),
\label{eq:gain-reparam}
\end{equation}
where $\gamma_{s,j}\in\mathbb R$ is the residual \emph{log-gain correction}, and
\begin{equation}
N_{s,j}
=
\phi(\eta_{s,j}),
\qquad
\phi(\eta)=\log(1+e^\eta),
\label{eq:noise-reparam}
\end{equation}
where $\eta_{s,j}\in\mathbb R$ is an unconstrained \emph{noise parameter}. Thus positivity of the additive floor is enforced by the softplus map \eqref{eq:noise-reparam}, not by post-hoc clipping.

The derivative of \eqref{eq:noise-reparam} is
\begin{equation}
\phi'(\eta)
=
\frac{1}{1+e^{-\eta}},
\label{eq:sigmoid}
\end{equation}
which is used in the Gauss--Newton linearization.

\subsection{Identifiability}

The model is invariant under
\[
G_{s,j}\mapsto c_j G_{s,j},
\qquad
S_{k,j}\mapsto c_j^{-1}S_{k,j},
\qquad c_j>0,
\]
so it is identifiable only up to a frequency-wise multiplicative scale.

The implementation resolves this by projecting the residual log-gains to zero mean across sensors:
\begin{equation}
\frac{1}{M}\sum_{s=1}^{M}\gamma_{s,j}=0,
\qquad j=1,\dots,J.
\label{eq:ident}
\end{equation}
Equivalently,
\begin{equation}
\frac{1}{M}\sum_{s=1}^{M}\log G_{s,j}
=
\frac{1}{M}\sum_{s=1}^{M}\log G^{(0)}_{s,j}
+
\frac{1}{M}\sum_{s=1}^{M}\gamma_{s,j}
=
0,
\label{eq:ident-full}
\end{equation}
because the active campaign workflow uses the unity baseline
\eqref{eq:nominal-baseline}. More generally, the core solver would only require
the baseline gains to have zero-mean log scale across sensors.

Under \eqref{eq:ident}, the latent spectra $\{S_{k,j}\}$ live on the common network reference scale.

\subsection{Reliable-sensor prior}

Let $s_\star$ denote the sensor believed to be the most reliable one. In the
current campaign workflow, this sensor is not chosen manually. Instead the code
first ranks sensors by record-wise cumulative cross-correlation, then filters
those candidates against a campaign-wide PSD-distribution similarity diagnostic.
The reliable sensor is the first record-wise winner that is not flagged as a
distribution-shape outlier; if all retained sensors are flagged, the code uses
the record-wise winner anyway.

Once $s_\star$ has been selected, the implementation does \emph{not} impose a
hard reference constraint. Instead it uses two soft mechanisms:

\begin{enumerate}
    \item A quadratic anchor on the residual log-gain of the reliable sensor, favoring
    \[
    \gamma_{s_\star,j}\approx 0,
    \]
    meaning that the reliable sensor should stay close to the unit network
    reference gain, or more generally to the nominal baseline supplied to the
    core solver.
    \item An optional multiplicative boost of its inverse-variance weight in the latent-spectrum update.
\end{enumerate}

This keeps the reliable sensor influential without forcing all mismatch into the rest of the network.

\subsection{Initialization}

Let $\mathcal K_{\mathrm{train}}\subseteq\{1,\dots,K\}$ be the training experiments used for fitting. The implementation initializes:
\begin{align}
\gamma^{(0)}_{s,j} &= 0,
\label{eq:init-gamma}
\\
G^{(0,\mathrm{fit})}_{s,j}
&=
G^{(0)}_{s,j},
\label{eq:init-gain}
\\
N^{(0)}_{s,j}
&=
0.1\,
Q_{0.05}\!\left(\{Y_{s,k,j}:k\in\mathcal K_{\mathrm{train}}\}\right),
\label{eq:init-noise}
\\
\eta^{(0)}_{s,j}
&=
\phi^{-1}(N^{(0)}_{s,j}),
\label{eq:init-eta}
\\
S^{(0)}_{k,j}
&=
\operatorname{median}_{s}
\left[
\frac{Y_{s,k,j}-N^{(0)}_{s,j}}{G^{(0,\mathrm{fit})}_{s,j}}
\right]_{+},
\label{eq:init-latent}
\end{align}
where $[\cdot]_+$ denotes clipping to a small positive floor for numerical stability.

The initial residual variance is
\begin{equation}
\sigma^{2(0)}_{s,j}
=
\max\!\left(
\frac{1}{|\mathcal K_{\mathrm{train}}|}
\sum_{k\in\mathcal K_{\mathrm{train}}}
\bigl(
Y_{s,k,j}-G^{(0,\mathrm{fit})}_{s,j}S^{(0)}_{k,j}-N^{(0)}_{s,j}
\bigr)^2,
\;
\sigma^2_{\min}
\right).
\label{eq:init-variance}
\end{equation}

\subsection{Second-difference smoothers}

Let $f_1<\dots<f_J$ denote the actual frequency grid used by the aligned
campaign tensor. The implementation does \emph{not} assume exact unit spacing.
Instead it normalizes the grid by the median spacing
\[
\tilde h = \operatorname{median}(f_{j+1}-f_j),
\qquad j=1,\dots,J-1,
\]
and defines normalized coordinates
\[
x_1=0,
\qquad
x_j=\sum_{u=1}^{j-1}\frac{f_{u+1}-f_u}{\tilde h},
\qquad j=2,\dots,J.
\]

For row $r\in\{1,\dots,J-2\}$, let
\[
h^-_r = x_{r+1}-x_r,
\qquad
h^+_r = x_{r+2}-x_{r+1}.
\]
The implemented curvature operator uses the three-point non-uniform second
difference
\[
(D_2 x)_r
=
\frac{2}{h^-_r(h^-_r+h^+_r)}x_r
-
\frac{2}{h^-_r h^+_r}x_{r+1}
+
\frac{2}{h^+_r(h^-_r+h^+_r)}x_{r+2}.
\]
The smoother then uses the corresponding Gram matrix
\begin{equation}
L = D_2^\top D_2.
\label{eq:L}
\end{equation}

Crucially, the gain smoother acts on $\gamma_s=(\gamma_{s,1},\dots,\gamma_{s,J})^\top$, i.e. on \emph{log-gain correction}, and the noise smoother acts on $\eta_s=(\eta_{s,1},\dots,\eta_{s,J})^\top$, i.e. on the \emph{softplus preimage}, not directly on $N_s$.

On an exactly uniform grid this reduces, up to the common scaling induced by
the normalized coordinates, to the usual second-difference smoother.

\subsection{Alternating estimation actually implemented}

The estimator runs for a fixed number of alternating iterations $t=0,\dots,T-1$.

\subsubsection*{Step 1: inverse-variance weights}

At iteration $t$, define
\begin{equation}
w^{(t)}_{s,j}
=
\frac{1}{\max(\sigma^{2(t)}_{s,j},\sigma^2_{\min})}.
\label{eq:base-weights}
\end{equation}
If a reliable-weight boost $\beta_{\mathrm{rel}}>0$ is used, then
\begin{equation}
w^{(t)}_{s_\star,j}
\leftarrow
\beta_{\mathrm{rel}}\,w^{(t)}_{s_\star,j}.
\label{eq:reliable-boost}
\end{equation}

\subsubsection*{Step 2: latent-spectrum update}

Given current gains and additive floors, the latent spectrum is updated by weighted least squares:
\begin{equation}
S^{(t+1)}_{k,j}
=
\left[
\frac{
\sum_{s=1}^{M}
w^{(t)}_{s,j}\,
G^{(t)}_{s,j}\,
\bigl(Y_{s,k,j}-N^{(t)}_{s,j}\bigr)
}{
\sum_{s=1}^{M}
w^{(t)}_{s,j}\,
\bigl(G^{(t)}_{s,j}\bigr)^2
}
\right]_{+}.
\label{eq:latent-update}
\end{equation}

\subsubsection*{Step 3: frequency-only low-information diagnostic}

The implementation measures how much the latent spectrum varies across calibration experiments:
\begin{equation}
v^{(t+1)}_j
=
\operatorname{Var}_{k\in\mathcal K_{\mathrm{train}}}
\bigl(S^{(t+1)}_{k,j}\bigr).
\label{eq:latent-var}
\end{equation}

Let
\[
\mathcal J_+ = \{j : v^{(t+1)}_j > \varepsilon\}.
\]
If $\mathcal J_+=\varnothing$, then every frequency is declared low-information and
\begin{equation}
\omega^{\mathrm{freq}}_j = \omega_{\min},
\qquad
m^{\mathrm{freq}}_j = 1,
\label{eq:freq-degenerate}
\end{equation}
for all $j$.

Otherwise define the threshold
\begin{equation}
\tau_v
=
\max\!\left(
\rho_v\,
\operatorname{median}\{v^{(t+1)}_j : j\in\mathcal J_+\},
\;
\varepsilon
\right),
\label{eq:var-threshold}
\end{equation}
where $\rho_v>0$ is the configured threshold ratio. Then
\begin{equation}
\omega^{\mathrm{freq}}_j
=
\operatorname{clip}
\left(
\frac{v^{(t+1)}_j}{\tau_v},
\;
\omega_{\min},
\;
1
\right),
\qquad
m^{\mathrm{freq}}_j
=
\mathbf 1\{v^{(t+1)}_j<\tau_v\}.
\label{eq:freq-weight}
\end{equation}

\subsubsection*{Step 4: sensor-specific information weights}

The frequency-only diagnostic is then refined into a sensor-specific information weight. Let
\begin{equation}
\bar S^{(t+1)}_j
=
\frac{1}{|\mathcal K_{\mathrm{train}}|}
\sum_{k\in\mathcal K_{\mathrm{train}}}
S^{(t+1)}_{k,j}.
\label{eq:latent-mean}
\end{equation}

Define the centered signal energy
\begin{equation}
E^{\mathrm{sig}}_{s,j}
=
\bigl(G^{(t)}_{s,j}\bigr)^2
\sum_{k\in\mathcal K_{\mathrm{train}}}
\bigl(S^{(t+1)}_{k,j}-\bar S^{(t+1)}_j\bigr)^2,
\label{eq:centered-energy}
\end{equation}
and the intercept-sensitive noise energy
\begin{equation}
E^{\mathrm{int}}_{s,j}
=
|\mathcal K_{\mathrm{train}}|\,
\bigl(\phi'(\eta^{(t)}_{s,j})\bigr)^2.
\label{eq:intercept-energy}
\end{equation}

The local design-balance score is
\begin{equation}
d_{s,j}
=
\frac{
\min\{E^{\mathrm{sig}}_{s,j},E^{\mathrm{int}}_{s,j}\}
}{
\max\{E^{\mathrm{sig}}_{s,j},E^{\mathrm{int}}_{s,j}\}+\varepsilon
}.
\label{eq:design-balance}
\end{equation}

The local residual-SNR score is
\begin{equation}
q_{s,j}
=
\frac{
\bigl(G^{(t)}_{s,j}\bigr)^2 v^{(t+1)}_j
}{
\bigl(G^{(t)}_{s,j}\bigr)^2 v^{(t+1)}_j + \sigma^{2(t)}_{s,j} + \varepsilon
}.
\label{eq:residual-snr}
\end{equation}

These are normalized by robust positive-entry medians:
\begin{align}
\tau_d
&=
\max\!\left(
\operatorname{median}\{d_{s,j}:d_{s,j}>0\},
\varepsilon
\right),
\label{eq:tau-d}
\\
\tau_q
&=
\max\!\left(
\operatorname{median}\{q_{s,j}:q_{s,j}>0\},
\varepsilon
\right),
\label{eq:tau-q}
\\
\widetilde d_{s,j}
&=
\operatorname{clip}\!\left(\frac{d_{s,j}}{\tau_d},0,1\right),
\label{eq:d-norm}
\\
\widetilde q_{s,j}
&=
\operatorname{clip}\!\left(\frac{q_{s,j}}{\tau_q},0,1\right).
\label{eq:q-norm}
\end{align}

The final sensor-specific information weight is
\begin{equation}
\omega_{s,j}
=
\operatorname{clip}
\bigl(
\omega^{\mathrm{freq}}_j\,
\widetilde d_{s,j}\,
\widetilde q_{s,j},
\;
\omega_{\min},
\;
1
\bigr).
\label{eq:information-weight}
\end{equation}

The implementation also stores the conservative warning mask
\begin{equation}
m_{s,j}
=
m^{\mathrm{freq}}_j
\;\vee\;
\mathbf 1\!\left\{
\omega_{s,j}\le \max(\omega_{\min},0.60)
\right\}.
\label{eq:low-info-mask}
\end{equation}

\subsubsection*{Step 5: one joint sensor update by sparse Gauss--Newton}

Fix a sensor $s$ and suppress the superscript $(t)$ for readability. Let
\begin{align}
g^{\mathrm{ref}}_j
&=
G^{(0)}_{s,j}\exp(\gamma_{s,j}),
\label{eq:gref}
\\
n^{\mathrm{ref}}_j
&=
\phi(\eta_{s,j}),
\label{eq:nref}
\\
h_j
&=
\phi'(\eta_{s,j}),
\label{eq:href}
\\
x_{k,j}
&=
g^{\mathrm{ref}}_j\,S_{k,j}.
\label{eq:xkj}
\end{align}

The nonlinear model is linearized around the current iterate:
\begin{equation}
Y_{s,k,j}
\approx
x_{k,j}
+
n^{\mathrm{ref}}_j
+
x_{k,j}\,\delta\gamma_j
+
h_j\,\delta\eta_j.
\label{eq:linearization}
\end{equation}

The effective data weight for that sensor and frequency is
\begin{equation}
\widetilde w_{s,j}
=
w_{s,j}\,\omega_{s,j}.
\label{eq:effective-weight}
\end{equation}

Let $\delta\gamma,\delta\eta\in\mathbb R^J$ collect the per-frequency increments. Define
\begin{align}
D^\gamma_s
&=
\operatorname{diag}\!\left(
\widetilde w_{s,j}\sum_{k}x_{k,j}^2
\right)_{j=1}^{J},
\label{eq:Dgamma}
\\
D^\eta_s
&=
\operatorname{diag}\!\left(
\widetilde w_{s,j}\,|\mathcal K_{\mathrm{train}}|\,h_j^2
\right)_{j=1}^{J},
\label{eq:Deta}
\\
C_s
&=
\operatorname{diag}\!\left(
\widetilde w_{s,j}\,h_j\sum_{k}x_{k,j}
\right)_{j=1}^{J},
\label{eq:Cs}
\\
b^\gamma_{s,j}
&=
\widetilde w_{s,j}\sum_{k}
x_{k,j}
\bigl(
Y_{s,k,j}-x_{k,j}-n^{\mathrm{ref}}_j
\bigr),
\label{eq:bgamma}
\\
b^\eta_{s,j}
&=
\widetilde w_{s,j}\,h_j\sum_{k}
\bigl(
Y_{s,k,j}-x_{k,j}-n^{\mathrm{ref}}_j
\bigr).
\label{eq:beta}
\end{align}

Let $\lambda_{G,\mathrm{ref}}$ and $\lambda_{N,\mathrm{ref}}$ be the conservative reference penalties. Normally the reference targets are the current iterates:
\[
\gamma^{\mathrm{tar}}_s=\gamma_s,
\qquad
\eta^{\mathrm{tar}}_s=\eta_s.
\]
For the reliable sensor $s_\star$, the gain target is blended with zero:
\begin{equation}
\gamma^{\mathrm{tar}}_{s_\star}
=
\frac{
\lambda_{G,\mathrm{ref}}\gamma_{s_\star}
+
\lambda_{\mathrm{rel}}\,0
}{
\lambda_{G,\mathrm{ref}}+\lambda_{\mathrm{rel}}
},
\qquad
\lambda^{(\star)}_{G,\mathrm{ref}}
=
\lambda_{G,\mathrm{ref}}+\lambda_{\mathrm{rel}}.
\label{eq:reliable-target}
\end{equation}
For $s\neq s_\star$, one simply has $\lambda^{(s)}_{G,\mathrm{ref}}=\lambda_{G,\mathrm{ref}}$.

The update solves the sparse block system
\begin{equation}
\resizebox{\textwidth}{!}{$
\begin{bmatrix}
D^\gamma_s + \lambda_{G,\mathrm{smooth}}L + \lambda^{(s)}_{G,\mathrm{ref}}I
&
C_s
\\[0.5em]
C_s
&
D^\eta_s + \lambda_{N,\mathrm{smooth}}L + \lambda_{N,\mathrm{ref}}I
\end{bmatrix}
\begin{bmatrix}
\delta\gamma_s\\[0.25em]
\delta\eta_s
\end{bmatrix}
=
\begin{bmatrix}
b^\gamma_s
-
\lambda_{G,\mathrm{smooth}}L\gamma_s
+
\lambda^{(s)}_{G,\mathrm{ref}}(\gamma^{\mathrm{tar}}_s-\gamma_s)
\\[0.5em]
b^\eta_s
-
\lambda_{N,\mathrm{smooth}}L\eta_s
+
\lambda_{N,\mathrm{ref}}(\eta^{\mathrm{tar}}_s-\eta_s)
\end{bmatrix}.
$}
\label{eq:block-system}
\end{equation}

The raw update is then
\begin{equation}
\gamma^{\mathrm{raw}}_s=\gamma_s+\delta\gamma_s,
\qquad
\eta^{\mathrm{raw}}_s=\eta_s+\delta\eta_s.
\label{eq:raw-update}
\end{equation}

If a correction cap $c_{\max}>0$ dB is configured, the corresponding log-gain bound is
\begin{equation}
\gamma_{\max}
=
\log\!\left(10^{c_{\max}/10}\right),
\label{eq:gamma-max}
\end{equation}
and each raw gain correction is clipped to
\begin{equation}
-\gamma_{\max}\le \gamma^{\mathrm{raw}}_{s,j}\le \gamma_{\max}.
\label{eq:raw-clip}
\end{equation}

\subsubsection*{Step 6: projection onto the identifiability constraint}

The sensor updates are computed independently, so their raw log-gain corrections generally do not satisfy \eqref{eq:ident}. The implementation therefore projects them frequency-by-frequency.

If no cap is active, then the projection is simply
\begin{equation}
\gamma^{(t+1)}_{s,j}
=
\gamma^{\mathrm{raw}}_{s,j}
-
\frac{1}{M}\sum_{u=1}^{M}\gamma^{\mathrm{raw}}_{u,j}.
\label{eq:projection-no-cap}
\end{equation}

If the cap is active, then for each $j$ the code finds a scalar offset $\alpha_j$ by bisection such that
\begin{equation}
\frac{1}{M}\sum_{s=1}^{M}
\operatorname{clip}
\bigl(
\gamma^{\mathrm{raw}}_{s,j}-\alpha_j,
-\gamma_{\max},
\gamma_{\max}
\bigr)
=
0,
\label{eq:projection-cap}
\end{equation}
and sets
\begin{equation}
\gamma^{(t+1)}_{s,j}
=
\operatorname{clip}
\bigl(
\gamma^{\mathrm{raw}}_{s,j}-\alpha_j,
-\gamma_{\max},
\gamma_{\max}
\bigr).
\label{eq:projection-cap-result}
\end{equation}

The frequency-wise common log shift removed by this projection is
\begin{equation}
\Delta_j
=
\frac{1}{M}\sum_{s=1}^{M}
\bigl(
\gamma^{\mathrm{raw}}_{s,j}-\gamma^{(t+1)}_{s,j}
\bigr),
\label{eq:common-shift}
\end{equation}
and is transferred into the latent spectra:
\begin{equation}
S^{(t+1)}_{k,j}
\leftarrow
S^{(t+1)}_{k,j}\exp(\Delta_j).
\label{eq:latent-rescale}
\end{equation}

Finally,
\begin{align}
G^{(t+1)}_{s,j}
&=
G^{(0)}_{s,j}\exp(\gamma^{(t+1)}_{s,j}),
\label{eq:g-update}
\\
N^{(t+1)}_{s,j}
&=
\phi(\eta^{\mathrm{raw}}_{s,j}).
\label{eq:n-update}
\end{align}

\subsubsection*{Step 7: residual-variance update}

With the new parameters,
\begin{equation}
r^{(t+1)}_{s,k,j}
=
Y_{s,k,j}
-
G^{(t+1)}_{s,j}S^{(t+1)}_{k,j}
-
N^{(t+1)}_{s,j},
\label{eq:residual}
\end{equation}
and
\begin{equation}
\sigma^{2(t+1)}_{s,j}
=
\max\!\left(
\frac{1}{|\mathcal K_{\mathrm{train}}|}
\sum_{k\in\mathcal K_{\mathrm{train}}}
\bigl(r^{(t+1)}_{s,k,j}\bigr)^2,
\;
\sigma^2_{\min}
\right).
\label{eq:variance-update}
\end{equation}

\subsection{Monitored objective}

The implementation stores, after each iteration, the diagnostic quantity
\begin{align}
\mathcal J_{\mathrm{diag}}
&=
\frac{1}{2}
\sum_{s=1}^{M}
\sum_{j=1}^{J}
\sum_{k\in\mathcal K_{\mathrm{train}}}
\left[
\log \sigma^2_{s,j}
+
\frac{r_{s,k,j}^2}{\sigma^2_{s,j}}
\right]
\nonumber\\
&\quad
+
\lambda_{G,\mathrm{smooth}}
\sum_{s=1}^{M}\|D_2\gamma_s\|_2^2
+
\lambda_{N,\mathrm{smooth}}
\sum_{s=1}^{M}\|D_2\eta_s\|_2^2
+
\lambda_{\mathrm{rel}}\|\gamma_{s_\star}\|_2^2.
\label{eq:diagnostic-objective}
\end{align}

This is the quantity recorded in \texttt{objective\_history}. It is a useful convergence diagnostic, but it is \emph{not} the full time-varying function minimized by every inner linear solve, because the actual update equations also include iteration-dependent conservative reference terms and information weights.

\subsection{Deployment-time calibration map}

For any new sensor-first tensor of observations $Y^{\mathrm{new}}_{s,\cdots,j}$, the deployed calibration map is
\begin{equation}
\widehat S^{\mathrm{new}}_{s,\cdots,j}
=
\frac{
Y^{\mathrm{new}}_{s,\cdots,j}-N_{s,j}
}{
G_{s,j}
}.
\label{eq:deploy}
\end{equation}
If nonnegativity is enforced, the implementation uses
\begin{equation}
\widehat S^{\mathrm{new}}_{s,\cdots,j}
=
\max\!\left\{
0,\,
\frac{
Y^{\mathrm{new}}_{s,\cdots,j}-N_{s,j}
}{
G_{s,j}
}
\right\}.
\label{eq:deploy-nonnegative}
\end{equation}

The code broadcasts the curves $(G_{s,j},N_{s,j})$ across any intermediate deployment dimensions while keeping the first axis as sensor and the last axis as frequency.

\subsection{Consensus estimator}

If several calibrated sensors later observe the same latent spectrum again, the implementation forms the inverse-variance consensus
\begin{equation}
\widehat S^{\mathrm{net}}_{k,j}
=
\frac{
\sum_{s=1}^{M}
\sigma^{-2}_{s,j}\,
\widehat S_{s,k,j}
}{
\sum_{s=1}^{M}
\sigma^{-2}_{s,j}
}.
\label{eq:consensus}
\end{equation}
This aggregation is meaningful only under a common-field assumption at deployment time.

\subsection{Hold-out validation and returned diagnostics}

Validation is performed by splitting whole calibration experiments into training and test subsets,
\[
\mathcal K_{\mathrm{train}}\cup\mathcal K_{\mathrm{test}}
=
\{1,\dots,K\},
\qquad
\mathcal K_{\mathrm{train}}\cap\mathcal K_{\mathrm{test}}=\varnothing,
\]
estimating the calibration on $\mathcal K_{\mathrm{train}}$, and checking whether the deployed correction \eqref{eq:deploy} reduces inter-sensor dispersion on $\mathcal K_{\mathrm{test}}$.

In the current campaign adapter, hold-out reporting also includes per-sensor
bias and RMSE to a leave-one-sensor-out corrected consensus, the fraction of
bins where $Y_{s,k,j}\le N_{s,j}$ would make the deployed correction invalid
before nonnegativity truncation, the median information weight, the fraction of
low-information bins, the fraction of bins that hit the gain cap, and the
median alignment timing error.

The fitted result stores, in addition to $(G,N,\sigma^2,S)$, the following diagnostics:
\begin{itemize}
    \item the residual correction gain $\exp(\gamma_{s,j})$;
    \item the latent-variation curve $v_j$;
    \item the frequency-only weights $\omega^{\mathrm{freq}}_j$ and mask $m^{\mathrm{freq}}_j$;
    \item the sensor-specific information weights $\omega_{s,j}$ and mask $m_{s,j}$;
    \item a mask of gain-correction bins that hit the configured cap;
    \item a mask of additive-noise bins numerically indistinguishable from zero.
\end{itemize}

\subsection{Summary}

The implemented estimator is therefore an alternating latent-variable scheme with the following specific structure:

\begin{enumerate}
    \item load one campaign dataset, align rows by timestamp, and select a reliable sensor from ranking diagnostics;
    \item set the nominal gain baseline to one;
    \item represent gain by a residual log correction and noise by a softplus parameter;
    \item update the latent spectra by weighted least squares;
    \item detect weakly identifiable frequencies from latent variation and refine this into sensor-specific information weights;
    \item update each sensor by one sparse, frequency-coupled Gauss--Newton solve with smoothness and conservative reference penalties;
    \item project the raw gain corrections onto the mean-zero identifiability constraint, with an optional symmetric box constraint in dB;
    \item update residual variances and repeat.
\end{enumerate}

This is the actual solution implemented in the code. It is deliberately more conservative than the earlier idealized description: in weak-information bins it leans on regularization, damping, and projection rather than pretending that gain and additive floor are cleanly identifiable from the data.

\end{document}
